\SourceLeaf{456}%
\noindent which, after using the values denoted by $\mathfrak{E},\mathfrak{F},\mathfrak{G}$ and $\mathfrak{E}',\mathfrak{F}',\mathfrak{G}'$, lead to the following equation:
\[
\left(\frac{\partial u}{\partial\zeta'}\right)^2=\left(\frac{\partial u'}{\partial\zeta}\right)^2.
\]

\noindent Accordingly, the linear differential expression
\[
ud\zeta\pm u'd\zeta'=d\varphi
\]
\noindent represents the total differential of a function $\varphi$ of the variables $\zeta,\zeta'$. This property can, however, hold for any two surfaces $S$ and $S'$ that can be developed onto one another only if the lower sign is taken. With the opposite assumption, when $S$ coincides with $S'$, the product $2ud\zeta$ would have to be a total differential for every surface, a circumstance that occurs only for developable surfaces.

The foregoing equation contains the starting point chosen by \textsc{Gauss} in No.\ 2 for the general solution of the problem of deformation.

This equation holds for any two surfaces $S$ and $S'$ that can be developed onto one another, but \textls[140]{not only} for such surfaces.

How \textsc{Gauss} intended to use it for this general solution is not immediately clear from No.\ 4 of the note. This question calls for further detailed investigation of the properties of the function $\varphi$ to be chosen in the particular case.
\begin{flushright}\textsc{Weingarten}.\end{flushright}

